\documentclass[11pt,a4paper]{article}
\usepackage[utf8]{inputenc}
\usepackage[T1]{fontenc}
\usepackage[french]{babel}
\usepackage{lmodern}
\usepackage{geometry}
\usepackage{graphicx}
\usepackage{booktabs}
\usepackage{tabularx}
\usepackage{float}
\usepackage{xcolor}
\usepackage[hidelinks]{hyperref}
\geometry{margin=2cm}
\setlength{\parskip}{0.7em}
\setlength{\parindent}{0pt}
\begin{document}
\title{Rapport analytique - KAZIMIRIUM 1.json}
\author{GitHub Copilot}
\date{2026-04-30T02:09:48.962844+00:00}
\maketitle
\tableofcontents
\newpage
\section*{Resume executif}
\addcontentsline{toc}{section}{Resume executif}
\begin{itemize}
\item \textbf{La partie presente un volume d'action eleve} (forte) : Le match cumule 535 mouvements, 49 constructions, 64 captures et 4 upgrades sur 138 tours.
\item \textbf{L'eau devient une vraie taxe de mobilite apres l'ouverture} (forte) : La perte moyenne passe de 0.04 destination par record en Q1 a 19.48 en Q3, soit 6.4\% de la mobilite seche disponible.
\item \textbf{Quand le royaume actif est couvert, le pilotage ralentit nettement} (forte) : Quand au moins une piece ennemie est cachee au royaume actif, les interactions montent de +7.53, les audits de +0.44, et le commit moyen prend +20.04 s de plus.
\item \textbf{Les coffres favorisent surtout des bonus de capacite sur cette partie} (moderee) : Le type de recompense le plus observe est build\_points\_max\_bonus, avec 3 occurrences.
\item \textbf{La boucle infernale s'ancre surtout sur une dette de camp, pas sur un total abstrait} (forte) : La dette moyenne du camp noir atteint 6.02 contre 3.22 pour l'autre camp, et 4 spawns infernaux sur 5 ciblent ce meme camp.
\item \textbf{Les fronts meteo reconfigurent fortement le pilotage} (forte) : Quand un front meteo est actif, la charge de decision change nettement: interactions +16.97, audits +2.00, avec +0.65 mouvements et +0.13 constructions par record.
\item \textbf{Les ouvertures de coffre se convertissent surtout en execution} (moderee) : Autour de une ouverture de coffre, l'effet visible porte d'abord sur le pilotage: interactions +4.96, audits +1.29, pour un effet de plateau de +0.42 mouvements et +0.21 constructions.
\item \textbf{Certaines conclusions doivent rester prudentes} (forte) : Les comparaisons de fenetres et de regimes sont utiles pour produire des hypotheses solides, mais elles ne valent pas preuve causale a elles seules.
\end{itemize}
\section{Methodologie}
Ce rapport repose sur le JSON d'analyse intermediaire produit a partir du companion schema-v5. Toutes les jointures temporelles utilisent le sequence index interne, et non le numero de tour commite, car ce dernier n'est pas suffisamment unique dans cet echantillon.
\section{Vue d'ensemble du match}
Cette partie oppose Max contre Ilan sur 138 tours, avec un volume d'action eleve et une telemetrie exploitable.
\subsection*{Indicateurs clefs}
\begin{tabularx}{\linewidth}{lX}
\toprule
Indicateur & Valeur \\
\midrule
Tours & 138 \\
Records & 276 \\
Mouvements & 535 \\
Constructions & 49 \\
Captures & 64 \\
Upgrades & 4 \\
\bottomrule
\end{tabularx}
\subsection*{La partie presente un volume d'action eleve}
\textbf{Niveau de confiance :} forte\\
Le match cumule 535 mouvements, 49 constructions, 64 captures et 4 upgrades sur 138 tours.
\begin{itemize}
\item Moyenne de mouvements par tour: 3.88.
\item Moyenne de constructions par tour: 0.36.
\item Moyenne d'audits de commandes par record: 3.36.
\end{itemize}
\subsection*{La telemetrie est suffisamment riche pour une lecture comportementale}
\textbf{Niveau de confiance :} forte\\
La partie dispose de 5828 evenements telemetry et d'une duree moyenne de record de 36939.14 ms.
\begin{itemize}
\item Behavioral telemetry active: oui.
\item Data collection active: oui.
\end{itemize}
\section{Terrain aleatoire et taxe de mobilite}
Le terrain genere aleatoirement n'est pas un simple decor: ici, l'eau retire durablement des destinations de mouvement et penalise surtout les pieces qui vivent sur des lignes ouvertes.
\subsection*{Indicateurs clefs}
\begin{tabularx}{\linewidth}{lX}
\toprule
Indicateur & Valeur \\
\midrule
Cellules eau & 70 \\
Carte jouable en eau & 3.5\% \\
Destinations perdues moyennes & 10.39 \\
Mobilite seche perdue & 4.0\% \\
\bottomrule
\end{tabularx}
\subsection*{L'eau devient une vraie taxe de mobilite apres l'ouverture}
\textbf{Niveau de confiance :} forte\\
La perte moyenne passe de 0.04 destination par record en Q1 a 19.48 en Q3, soit 6.4\% de la mobilite seche disponible.
\begin{itemize}
\item Couverture eau de la carte jouable: 3.5\%.
\item Perte moyenne globale: 10.39 destinations par record actif.
\item Nombre moyen de pieces impactees: 2.10.
\end{itemize}
\textit{Le calcul compare chaque snapshot a une version seche du meme plateau; quelques etats speciaux de breche de mur restent approximes.}
\subsection*{La penalite touche surtout les pieces de longue portee}
\textbf{Niveau de confiance :} forte\\
Les fous, tours et dames concentrent 99.9\% des destinations perdues a cause de l'eau, alors que les pions n'en perdent aucune sur cet echantillon.
\begin{itemize}
\item Piece la plus penalisee en part relative: Queen (10.3\%).
\item Fous: 6.2\% de mobilite seche perdue.
\item Tours: 8.4\% de mobilite seche perdue.
\item Dames: 10.3\% de mobilite seche perdue.
\end{itemize}
\begin{figure}[H]
\centering
\includegraphics[width=0.98\linewidth]{\detokenize{charts/terrain_water_tax.png}}
\caption{Evolution de la mobilite seche perdue a cause de l'eau au fil du match, et sensibilite par type de piece.}
\end{figure}
\subsection*{Taxe eau par quartile de match}
\begin{tabularx}{\linewidth}{X|X|X|X|X}
\toprule
Quartile & Records & Destinations perdues moyennes & Part de mobilite seche perdue & Pieces impactees moyennes \\
\midrule
Q1 & 69 & 0.04 & 0.3\% & 0.03 \\
Q2 & 69 & 11.74 & 5.5\% & 2.29 \\
Q3 & 69 & 19.48 & 6.4\% & 3.70 \\
Q4 & 69 & 10.29 & 3.7\% & 2.39 \\
\bottomrule
\end{tabularx}
\subsection*{Sensibilite a l'eau par type de piece}
\begin{tabularx}{\linewidth}{X|X|X|X|X}
\toprule
Piece & Echantillons actifs & Destinations perdues & Part de mobilite perdue & Taux d'impact \\
\midrule
Queen & 31 & 83 & 10.3\% & 25.8\% \\
Rook & 150 & 280 & 8.4\% & 43.3\% \\
Bishop & 1839 & 2501 & 6.2\% & 27.4\% \\
King & 276 & 3 & 0.1\% & 1.1\% \\
Pawn & 3334 & 0 & 0.0\% & 0.0\% \\
\bottomrule
\end{tabularx}
\subsection*{Taxe eau par royaume actif}
\begin{tabularx}{\linewidth}{X|X|X|X|X}
\toprule
Royaume & Records & Destinations perdues moyennes & Part de mobilite perdue & Pieces impactees moyennes \\
\midrule
black & 138 & 12.23 & 5.4\% & 2.61 \\
white & 138 & 8.54 & 2.6\% & 1.59 \\
\bottomrule
\end{tabularx}
\subsection*{Pics de pression eau}
\begin{tabularx}{\linewidth}{X|X|X|X|X}
\toprule
Tour & Royaume actif & Pieces actives & Destinations perdues & Part de mobilite perdue \\
\midrule
68 & white & 33 & 36 & 8.7\% \\
69 & white & 32 & 33 & 8.7\% \\
58 & black & 22 & 32 & 12.1\% \\
89 & black & 29 & 32 & 10.3\% \\
53 & black & 18 & 31 & 13.7\% \\
\bottomrule
\end{tabularx}
\section{Meteo et brouillard aleatoire}
La meteo est etudiee comme un systeme aleatoire de couverture: quelle surface est dissimulante, combien de pieces chaque camp perd reellement de vue, et comment ce brouillard change le pilotage du joueur actif.
\subsection*{Indicateurs clefs}
\begin{tabularx}{\linewidth}{lX}
\toprule
Indicateur & Valeur \\
\midrule
Records meteo & 276 \\
Records avec front actif & 67.0\% \\
Brouillard dissimulant moyen & 144.82 \\
Pic de pieces masquees & 26 \\
Records couverts & 99 \\
\bottomrule
\end{tabularx}
\subsection*{Quand le royaume actif est couvert, le pilotage ralentit nettement}
\textbf{Niveau de confiance :} forte\\
Quand au moins une piece ennemie est cachee au royaume actif, les interactions montent de +7.53, les audits de +0.44, et le commit moyen prend +20.04 s de plus.
\begin{itemize}
\item Records couverts: 99, records degages: 177.
\item Ennemis caches moyens: couvert=7.46, degage=0.
\item Brouillard dissimulant moyen: couvert=254.59, degage=83.42.
\end{itemize}
\textit{Lecture comparative: la couverture agit souvent avec d'autres contraintes de milieu, elle ne prouve pas a elle seule la causalite.}
\subsection*{Le brouillard ne masque pas les deux camps de la meme facon}
\textbf{Niveau de confiance :} forte\\
Sur cet echantillon, le camp noir joue en moyenne avec 2.97 pieces ennemies masquees contre 2.31 pour le camp blanc.
\begin{itemize}
\item Pic de pieces ennemies masquees: blanc=17, noir=26.
\item Records avec au moins une piece masquee: blanc=38.4\%, noir=33.7\%
\end{itemize}
\subsection*{La meteo devient structurelle en fin de partie}
\textbf{Niveau de confiance :} forte\\
Le front actif passe d'environ 0.0\% dans le premier quartile a 92.8\% dans le dernier, ce qui signale une presence beaucoup plus continue en fin de match.
\begin{itemize}
\item Taux d'activite global des fronts: 67.0\%.
\item Brouillard dissimulant moyen: 144.82 cellules.
\end{itemize}
\begin{figure}[H]
\centering
\includegraphics[width=0.98\linewidth]{\detokenize{charts/weather_pressure.png}}
\caption{En haut, la surface nuageuse totale et dissimulante; au milieu, les pieces cachees a chaque camp; en bas, l'ecart relatif de comportement quand le royaume actif perd reellement de la vision.}
\end{figure}
\subsection*{Regimes meteo par quartile de match}
\begin{tabularx}{\linewidth}{X|X|X|X|X}
\toprule
Quartile & Records & Front actif & Brouillard dissimulant moyen & Pieces ennemies masquees moyennes \\
\midrule
Q1 & 69 & 0.0\% & 0 & 0 \\
Q2 & 69 & 75.4\% & 216.84 & 5.90 \\
Q3 & 69 & 100.0\% & 154.38 & 5.84 \\
Q4 & 69 & 92.8\% & 208.06 & 9.38 \\
\bottomrule
\end{tabularx}
\subsection*{Visibilite masquee par camp}
\begin{tabularx}{\linewidth}{X|X|X|X}
\toprule
Camp & Pieces ennemies masquees moyennes & Pic masque & Part des records avec masque \\
\midrule
white & 2.31 & 17 & 38.4\% \\
black & 2.97 & 26 & 33.7\% \\
\bottomrule
\end{tabularx}
\subsection*{Comportement du royaume actif sous couverture}
\begin{tabularx}{\linewidth}{X|X|X|X}
\toprule
Metrique & Couvert & Degage & Delta couvert - degage \\
\midrule
Ennemis caches moyens & 7.46 & 0 & +7.46 \\
Brouillard dissimulant moyen & 254.59 & 83.42 & +171.16 \\
Moves moyens & 2.05 & 1.88 & +0.17 \\
Builds moyens & 0.16 & 0.19 & -0.02 \\
Audits moyens & 3.65 & 3.20 & +0.44 \\
Interactions moyennes & 23.44 & 15.92 & +7.53 \\
Commit moyen (s) & 49.79 & 29.75 & +20.04 \\
\bottomrule
\end{tabularx}
\section{Coffres et recompenses aleatoires}
Les coffres sont analyses comme un systeme de gratification aleatoire: quelle recompense sort, et a quel rythme les joueurs la convertissent en actions.
\subsection*{Indicateurs clefs}
\begin{tabularx}{\linewidth}{lX}
\toprule
Indicateur & Valeur \\
\midrule
Spawns de coffres & 9 \\
Ouvertures de coffres & 8 \\
Taux d'ouverture & 88.9\% \\
Delai max avant collecte (tours intermediaires) & 48 \\
\bottomrule
\end{tabularx}
\subsection*{Les coffres favorisent surtout des bonus de capacite sur cette partie}
\textbf{Niveau de confiance :} moderee\\
Le type de recompense le plus observe est build\_points\_max\_bonus, avec 3 occurrences.
\begin{itemize}
\item Part des recompenses de bonus (mouvement/build): 75.0\%.
\item Taux d'ouverture par spawn: 88.9\%.
\end{itemize}
\subsection*{Les coffres observes pesent presque uniquement en phase tardive}
\textbf{Niveau de confiance :} moderee\\
Sur cette partie, 8 ouvertures sur 8 arrivent apres le seuil tardif: l'effet aleatoire mesurable provient donc surtout du mix de fin de partie.
\begin{itemize}
\item Ouvertures early: 0.
\item Ouvertures late: 8.
\item Poids configures tardifs: buildBonus=37.5\%, gold=25.0\%, movementBonus=37.5\%.
\item Mix observe tardif: buildBonus=37.5\%, gold=25.0\%, movementBonus=37.5\%.
\end{itemize}
\textit{Aucune ouverture en phase early: cet echantillon ne permet pas de verifier la bascule early->late des coffres.}
\subsection*{La collecte des coffres alterne entre opportunisme rapide et attente longue}
\textbf{Niveau de confiance :} moderee\\
Le delai moyen avant collecte est de 10 tours intermediaires, avec un maximum observe de 48 tours intermediaires.
\begin{itemize}
\item Nombre de coffres correctement apparies spawn->ouverture: 8.
\item Le delai est mesure ici comme le nombre de tours complets entre le spawn et l'ouverture, hors tour du spawn et tour de collecte.
\end{itemize}
\begin{figure}[H]
\centering
\includegraphics[width=0.98\linewidth]{\detokenize{charts/chest_rewards.png}}
\caption{Distribution des recompenses et histogramme des delais de collecte des coffres.}
\end{figure}
\subsection*{Distribution des recompenses de coffres}
\begin{tabularx}{\linewidth}{X|X}
\toprule
Type de recompense & Occurrences \\
\midrule
build\_points\_max\_bonus & 3 \\
gold & 2 \\
movement\_points\_max\_bonus & 3 \\
\bottomrule
\end{tabularx}
\subsection*{Delais de collecte des coffres}
\begin{tabularx}{\linewidth}{X|X|X|X|X|X}
\toprule
Tour spawn & Tour ouverture & Ecart brut (tours) & Tours intermediaires & Delai (records) & Recompense \\
\midrule
10 & 20 & 10 & 9 & 19 & build\_points\_max\_bonus \\
23 & 31 & 8 & 7 & 16 & build\_points\_max\_bonus \\
37 & 42 & 5 & 4 & 10 & build\_points\_max\_bonus \\
50 & 55 & 5 & 4 & 10 & gold \\
61 & 110 & 49 & 48 & 97 & gold \\
113 & 116 & 3 & 2 & 5 & movement\_points\_max\_bonus \\
119 & 120 & 1 & 0 & 2 & movement\_points\_max\_bonus \\
124 & 131 & 7 & 6 & 14 & movement\_points\_max\_bonus \\
\bottomrule
\end{tabularx}
\section{Systeme infernal et dette de sang}
Le systeme infernal est lu comme une boucle de dette par camp: on suit separement la dette blanche et la dette noire, puis on relie les spawns a la cible, a la piece manifestee et au niveau de dette atteint au moment du declenchement.
\subsection*{Indicateurs clefs}
\begin{tabularx}{\linewidth}{lX}
\toprule
Indicateur & Valeur \\
\midrule
Spawns infernaux & 5 \\
Lifetime moyen & 14 \\
Dette blanche moyenne & 3.22 \\
Dette noire moyenne & 6.02 \\
\bottomrule
\end{tabularx}
\subsection*{La boucle infernale s'ancre surtout sur une dette de camp, pas sur un total abstrait}
\textbf{Niveau de confiance :} forte\\
La dette moyenne du camp noir atteint 6.02 contre 3.22 pour l'autre camp, et 4 spawns infernaux sur 5 ciblent ce meme camp.
\begin{itemize}
\item Pics de dette: blanc=23, noir=27.
\item Cibles de spawn: black=4, white=1.
\item Tours de spawn observes: 25, 75, 85, 95, 120.
\end{itemize}
\subsection*{La dette infernale semble comprimer la cadence de spawn}
\textbf{Niveau de confiance :} exploratoire\\
Quand la cible porte deja une dette au moment d'un spawn, l'intervalle moyen jusqu'au spawn suivant tombe a 15 tours, contre 50 lorsqu'il n'y a pas de dette cible.
\begin{itemize}
\item Nombre de spawns observes: 5.
\item Lifetime moyen observe: 14 tours.
\item La mesure porte ici sur la prochaine cadence observee apres le spawn, ce qui capture mieux l'effet de la dette cible sur la boucle infernale.
\end{itemize}
\textit{Lecture exploratoire: peu de transitions comparables pour isoler proprement l'effet de la dette.}
\begin{figure}[H]
\centering
\includegraphics[width=0.98\linewidth]{\detokenize{charts/infernal_pressure.png}}
\caption{En haut, les courbes de dette blanche/noire avec les spawns annotes par piece et camp cible; en bas, le niveau de dette de chaque camp au moment exact des manifestations.}
\end{figure}
\subsection*{Dette infernale par royaume}
\begin{tabularx}{\linewidth}{X|X|X}
\toprule
Camp & Dette moyenne & Pic de dette \\
\midrule
white & 3.22 & 23 \\
black & 6.02 & 27 \\
\bottomrule
\end{tabularx}
\subsection*{Dette observee au moment des spawns infernaux}
\begin{tabularx}{\linewidth}{X|X|X|X|X|X}
\toprule
Tour & Cible & Piece manifestee & Dette blanche & Dette noire & Intervalle jusqu'au spawn suivant \\
\midrule
25 & black & pawn & 0 & 0 & 50 \\
75 & white & bishop & 3 & 6 & 10 \\
85 & black & rook & 1 & 5 & 10 \\
95 & black & pawn & 3 & 11 & 25 \\
120 & black & pawn & 11 & 14 & n/a \\
\bottomrule
\end{tabularx}
\section{Comportement joueur et charge de decision}
La telemetrie de comportement permet de voir si les systemes aleatoires changent la densite d'actions et le temps de decision.
\subsection*{Indicateurs clefs}
\begin{tabularx}{\linewidth}{lX}
\toprule
Indicateur & Valeur \\
\midrule
Audit moyen par record & 3.36 \\
Interactions moyennes par record & 18.62 \\
Delai moyen premiere action (ms) & 6166.70 \\
Delai moyen commit (ms) & 36939.14 \\
\bottomrule
\end{tabularx}
\subsection*{Les fronts meteo reconfigurent fortement le pilotage}
\textbf{Niveau de confiance :} forte\\
Quand un front meteo est actif, la charge de decision change nettement: interactions +16.97, audits +2.00, avec +0.65 mouvements et +0.13 constructions par record.
\begin{itemize}
\item Records front actif: 185, records sans front: 91.
\item Interactions moyennes: front actif=24.21, sans front=7.24.
\item Audits moyens: front actif=4.02, sans front=2.02.
\item Moves/builds moyens: front actif=2.15/0.22, sans front=1.51/0.09.
\end{itemize}
\subsection*{La pression infernale change surtout la charge de decision}
\textbf{Niveau de confiance :} forte\\
Quand au moins un infernal est actif, la charge de decision change nettement: interactions +5.28, audits +1.06, avec +0.41 mouvements et +0.14 constructions par record.
\begin{itemize}
\item Records infernal actif: 133, records sans infernal: 143.
\item Interactions moyennes: infernal actif=21.35, sans infernal=16.07.
\item Audits moyens: infernal actif=3.91, sans infernal=2.85.
\item Moves/builds moyens: infernal actif=2.15/0.25, sans infernal=1.74/0.11.
\end{itemize}
\subsection*{Les joueurs reajustent regulierement leurs plans}
\textbf{Niveau de confiance :} forte\\
Le match contient 64 remplacements de commandes et 69 annulations, signe d'un pilotage iteratif plutot que lineaire.
\begin{itemize}
\item Delai moyen premiere action: 6166.70 ms.
\item Delai moyen commit: 36939.14 ms.
\end{itemize}
\begin{figure}[H]
\centering
\includegraphics[width=0.98\linewidth]{\detokenize{charts/behavior_regimes.png}}
\caption{Variation relative des metriques de comportement entre regimes aleatoires actifs et inactifs, avec annotation du delta absolu par record.}
\end{figure}
\subsection*{Comparaison de regimes meteo}
\begin{tabularx}{\linewidth}{X|X|X|X}
\toprule
Metrique & Front actif & Sans front & Delta actif - sans front \\
\midrule
Moves moyens & 2.15 & 1.51 & +0.65 \\
Builds moyens & 0.22 & 0.09 & +0.13 \\
Audits moyens & 4.02 & 2.02 & +2.00 \\
Interactions moyennes & 24.21 & 7.24 & +16.97 \\
\bottomrule
\end{tabularx}
\subsection*{Comparaison de regimes infernaux}
\begin{tabularx}{\linewidth}{X|X|X|X}
\toprule
Metrique & Infernal actif & Sans infernal & Delta actif - sans infernal \\
\midrule
Moves moyens & 2.15 & 1.74 & +0.41 \\
Builds moyens & 0.25 & 0.11 & +0.14 \\
Audits moyens & 3.91 & 2.85 & +1.06 \\
Interactions moyennes & 21.35 & 16.07 & +5.28 \\
\bottomrule
\end{tabularx}
\section{Effets croises entre aleatoire et comportement}
Cette section croise les fenetres avant/apres evenement pour voir si un systeme aleatoire modifie le tempo observable juste autour de son apparition.
\subsection*{Indicateurs clefs}
\begin{tabularx}{\linewidth}{lX}
\toprule
Indicateur & Valeur \\
\midrule
Fenetres meteo & 4 \\
Fenetres coffres & 8 \\
Fenetres infernales & 5 \\
\bottomrule
\end{tabularx}
\subsection*{Les ouvertures de coffre se convertissent surtout en execution}
\textbf{Niveau de confiance :} moderee\\
Autour de une ouverture de coffre, l'effet visible porte d'abord sur le pilotage: interactions +4.96, audits +1.29, pour un effet de plateau de +0.42 mouvements et +0.21 constructions.
\begin{itemize}
\item Fenetres observees: 8.
\item Delta mouvements: +0.42.
\item Delta constructions: +0.21.
\item Delta audits/interactions: +1.29 / +4.96.
\end{itemize}
\textit{Les fenetres restent correlationnelles: elles montrent ce qui se passe juste avant/apres, pas une causalite prouvee.}
\subsection*{Les spawns infernaux provoquent surtout un surcroit de pilotage}
\textbf{Niveau de confiance :} exploratoire\\
Autour de un spawn infernal, l'effet visible porte d'abord sur le pilotage: interactions +3, audits +0.13, pour un effet de plateau de +0.27 mouvements et +0.07 constructions.
\begin{itemize}
\item Fenetres observees: 5.
\item Delta mouvements: +0.27.
\item Delta constructions: +0.07.
\item Delta audits/interactions: +0.13 / +3.00.
\end{itemize}
\textit{Les fenetres restent correlationnelles: elles montrent ce qui se passe juste avant/apres, pas une causalite prouvee.}
\subsection*{Les fronts meteo declenchent d'abord de la replanification}
\textbf{Niveau de confiance :} exploratoire\\
Autour de un spawn de front meteo, l'effet visible porte d'abord sur le pilotage: interactions +2.50, audits +0.92, pour un effet de plateau de -0.33 mouvements et +0.08 constructions.
\begin{itemize}
\item Fenetres observees: 4.
\item Delta mouvements: -0.33.
\item Delta constructions: +0.08.
\item Delta audits/interactions: +0.92 / +2.50.
\end{itemize}
\textit{Les fenetres restent correlationnelles: elles montrent ce qui se passe juste avant/apres, pas une causalite prouvee.}
\begin{figure}[H]
\centering
\includegraphics[width=0.98\linewidth]{\detokenize{charts/event_windows.png}}
\caption{Chaque sous-graphe montre le delta moyen avant/apres evenement sur les metriques de comportement qui appuient directement le commentaire du rapport.}
\end{figure}
\subsection*{Deltas moyens apres evenement}
\begin{tabularx}{\linewidth}{X|X|X|X|X}
\toprule
Evenement & Delta moves & Delta builds & Delta audits & Delta interactions \\
\midrule
Spawn front meteo & -0.33 & +0.08 & +0.92 & +2.50 \\
Ouverture coffre & +0.42 & +0.21 & +1.29 & +4.96 \\
Spawn infernal & +0.27 & +0.07 & +0.13 & +3.00 \\
\bottomrule
\end{tabularx}
\section{Limites d'interpretation}
Le rapport doit rester credibile: les limites des donnees sont explicites et ne doivent pas etre masquees par les graphiques.
\subsection*{Indicateurs clefs}
\begin{tabularx}{\linewidth}{lX}
\toprule
Indicateur & Valeur \\
\midrule
Nombre de limitations & 2 \\
\bottomrule
\end{tabularx}
\subsection*{Certaines conclusions doivent rester prudentes}
\textbf{Niveau de confiance :} forte\\
Les comparaisons de fenetres et de regimes sont utiles pour produire des hypotheses solides, mais elles ne valent pas preuve causale a elles seules.
\begin{itemize}
\item Le JSON d'analyse declare explicitement des limites sur la causalite.
\item Nombre de limitations source remontees: 2.
\end{itemize}
\subsection*{Limitations source}
\begin{tabularx}{\linewidth}{X|X|X}
\toprule
Scope & Niveau & Message \\
\midrule
xp & direct & No XP audit entries were found in the current companion. \\
causality & proxy & Cross-system comparisons are emitted as correlation-ready windows, not as proven causal effects. \\
\bottomrule
\end{tabularx}
\end{document}
